\documentclass[11pt,a4paper]{article}

% -----------------------------------------------------------------------------
% PACKAGES & CONFIGURATION
% -----------------------------------------------------------------------------
\usepackage[utf8]{inputenc}
\usepackage[T1]{fontenc}
\usepackage{lmodern}
\usepackage{microtype}
\usepackage[margin=1.8cm,top=2.0cm,bottom=2.0cm,headheight=14pt]{geometry}
\usepackage{xcolor}
\usepackage{fancyhdr}
\usepackage{titlesec}
\usepackage{booktabs}
\usepackage{enumitem}
\usepackage{tcolorbox}
\tcbuselibrary{skins}
\usepackage{fancyvrb}
\usepackage{etoolbox}
\usepackage{upquote}
\usepackage{hyperref}

% -----------------------------------------------------------------------------
% COLOR DEFINITIONS (ENSIA Lab Theme)
% -----------------------------------------------------------------------------
\definecolor{primary}{RGB}{31, 78, 121}        % Deep Navy Blue (#1f4e79)
\definecolor{darktext}{RGB}{30, 30, 30}        % Text color
\definecolor{codebg}{RGB}{246, 247, 249}       % Shaded code box background
\definecolor{codeborder}{RGB}{195, 205, 215}   % Code box border
\definecolor{barcolor}{RGB}{31, 78, 121}       % Left accent bar
\definecolor{notebg}{RGB}{238, 244, 250}       % Soft note box background

% -----------------------------------------------------------------------------
% HEADERS & FOOTERS (Matching ENSIA Lab Sheet)
% -----------------------------------------------------------------------------
\pagestyle{fancy}
\fancyhf{}
\fancyhead[L]{\small\textit{ENSIA -- AI and Data Science}}
\fancyhead[R]{\small\textit{Introduction to Mobile Robotics}}
\fancyfoot[L]{\small\textit{macOS Setup Guide: ROS 2 Jazzy (RoboStack + Pixi) \& CoppeliaSim}}
\fancyfoot[R]{\small\textit{Page \thepage}}
\renewcommand{\headrulewidth}{0.5pt}
\renewcommand{\footrulewidth}{0.5pt}

% -----------------------------------------------------------------------------
% SECTION TITLE FORMATTING
% -----------------------------------------------------------------------------
\titleformat{\section}
  {\color{primary}\Large\bfseries}{\thesection.}{0.6em}{}

\titleformat{\subsection}
  {\color{primary}\large\bfseries}{}{0pt}{}

\titlespacing*{\section}{0pt}{10pt plus 1pt minus 1pt}{5pt}
\titlespacing*{\subsection}{0pt}{8pt plus 1pt minus 1pt}{4pt}

% -----------------------------------------------------------------------------
% BOX STYLING
% -----------------------------------------------------------------------------
\newtcolorbox{cmdtcb}{
  enhanced,
  colback=codebg,
  colframe=barcolor,
  arc=1.0mm,
  leftrule=2.5pt,
  rightrule=0.4pt,
  toprule=0.4pt,
  bottomrule=0.4pt,
  boxsep=0.2mm,
  left=2.5mm,
  right=2.5mm,
  top=0.5mm,
  bottom=0.5mm
}
\DefineVerbatimEnvironment{cmdbox}{Verbatim}{fontsize=\footnotesize}
\BeforeBeginEnvironment{cmdbox}{\begin{cmdtcb}}
\AfterEndEnvironment{cmdbox}{\end{cmdtcb}}

\newtcolorbox{notebox}[1][Note]{
  enhanced,
  colback=notebg,
  colframe=barcolor,
  arc=1.0mm,
  leftrule=2.5pt,
  rightrule=0.4pt,
  toprule=0.4pt,
  bottomrule=0.4pt,
  boxsep=1.0mm,
  left=2.5mm,
  right=2.5mm,
  top=1.5mm,
  bottom=1.5mm,
  title={\textbf{#1}},
  coltitle=primary,
  attach title to upper={\par\vspace{2pt}}
}

\hypersetup{
  colorlinks=true,
  linkcolor=primary,
  citecolor=primary,
  urlcolor=primary,
  pdftitle={macOS Setup Guide: ROS 2 Jazzy (RoboStack + Pixi) \& CoppeliaSim}
}

% -----------------------------------------------------------------------------
% DOCUMENT BODY
% -----------------------------------------------------------------------------
\begin{document}

\begin{center}
  {\small ENSIA -- Artificial Intelligence and Data Science \quad \textbullet \quad Semester 7}\\[0.2cm]
  {\LARGE \textbf{\color{primary}macOS Setup Guide: ROS 2 Jazzy \& CoppeliaSim}}\\[0.15cm]
  {\normalsize Native Apple Silicon (\texttt{arm64}) Installation via RoboStack + Pixi}\\[0.25cm]
  \rule{\linewidth}{0.8pt}
\end{center}
\vspace{-0.1cm}

This guide lets macOS students run \textbf{ROS 2 Jazzy} natively alongside \textbf{CoppeliaSim} without a virtual machine, and includes the macOS-specific build fix for the \texttt{libsimExtROS2.dylib} plugin crash.

% =============================================================================
\section{Install Pixi \& ROS 2 Jazzy (RoboStack)}
% =============================================================================

\begin{enumerate}[leftmargin=1.5em]
  \item Install \textbf{Pixi} (the package manager recommended by RoboStack):
\begin{cmdbox}
curl -fsSL https://pixi.sh/install.sh | bash
\end{cmdbox}
  Restart your terminal (or run \texttt{source \textasciitilde/.zshrc}) so \texttt{pixi} is on your \texttt{PATH}.

  \item Create a Pixi workspace with the RoboStack Jazzy channel:
\begin{cmdbox}
mkdir -p ~/ros_jazzy_pixi && cd ~/ros_jazzy_pixi
pixi init -c https://prefix.dev/robostack-jazzy -c https://prefix.dev/conda-forge
\end{cmdbox}

  \item Add ROS 2 Jazzy Desktop, build tools, and CoppeliaSim's Python IPC dependencies (\texttt{pyzmq}, \texttt{cbor2}):
\begin{cmdbox}
pixi add ros-jazzy-desktop colcon-common-extensions
pixi add ros-jazzy-teleop-twist-keyboard
pixi add ros-jazzy-plotjuggler-ros pyzmq cbor2
\end{cmdbox}

  \item \textit{(Recommended)} Pin local discovery in \texttt{\textasciitilde/ros\_jazzy\_pixi/pixi.toml} so nodes on the same Mac always discover each other cleanly:
\begin{cmdbox}
printf '\n[activation.env]\n' >> ~/ros_jazzy_pixi/pixi.toml
printf 'ROS_AUTOMATIC_DISCOVERY_RANGE = "LOCALHOST"\n' >> ~/ros_jazzy_pixi/pixi.toml
printf 'ROS_DOMAIN_ID = "42"\n' >> ~/ros_jazzy_pixi/pixi.toml
\end{cmdbox}

  \item Verify ROS 2 works by running the talker and listener in two terminals:
\begin{cmdbox}
cd ~/ros_jazzy_pixi
pixi run ros2 run demo_nodes_cpp talker
\end{cmdbox}
\begin{cmdbox}
cd ~/ros_jazzy_pixi
pixi run ros2 run demo_nodes_py listener
\end{cmdbox}
  You should see \texttt{Hello World} messages published and received.
\end{enumerate}

% =============================================================================
\section{Install CoppeliaSim}
% =============================================================================

\begin{enumerate}[leftmargin=1.5em]
  \item Download \textbf{CoppeliaSim Edu} for \textbf{macOS (Apple Silicon / arm64)} from the official downloads page:\\
  \url{https://downloads.coppeliarobotics.com} \quad (or \url{https://www.coppeliarobotics.com}).
  \item Open the downloaded archive and move \texttt{coppeliaSim.app} into \texttt{/Applications/coppeliaSim.app}.
\end{enumerate}

% =============================================================================
\section{Build \texttt{sim\_ros2\_interface} with the macOS Python Fix}
% =============================================================================

\begin{notebox}[Issue Description (macOS)]
When starting a ROS 2 simulation in CoppeliaSim on macOS, the simulation fails by default with:
\begin{center}
\texttt{[CoppeliaSim:error] Cannot load library libsimExtROS2.dylib: (dlopen(...): symbol not found in flat namespace '\_PyExc\_RuntimeError')} \quad and \quad \texttt{plugin 'ROS2': load failed.}
\end{center}
ROS 2 message packages (\texttt{rosidl\_generator\_py}) depend on Python C-API symbols, but the upstream \texttt{sim\_ros2\_interface} \texttt{CMakeLists.txt} only searches for \texttt{Interpreter} without linking \texttt{Python3::Python}. Under macOS \texttt{dyld}, unlinked Python symbols cannot be resolved at runtime.
\end{notebox}

\begin{enumerate}[leftmargin=1.5em]
  \item Clone \texttt{simROS2} into a ROS 2 workspace:
\begin{cmdbox}
mkdir -p ~/ros2_ws/src && cd ~/ros2_ws/src
git clone https://github.com/CoppeliaRobotics/simROS2.git sim_ros2_interface
\end{cmdbox}

  \item Patch \texttt{CMakeLists.txt} to explicitly link \texttt{Python3::Python}:
\begin{cmdbox}
python3 -c '
import os
p = os.path.expanduser("~/ros2_ws/src/sim_ros2_interface/CMakeLists.txt")
s = open(p).read()
s = s.replace("find_package(Python3 REQUIRED COMPONENTS Interpreter)",
              "find_package(Python3 REQUIRED COMPONENTS Interpreter Development)", 1)
s = s.replace("target_link_libraries(simROS2 PRIVATE Boost::boost)",
              "target_link_libraries(simROS2 PRIVATE Boost::boost Python3::Python)", 1)
open(p, "w").write(s)
'
\end{cmdbox}

  \item Add the lab message types (\texttt{Pose2D}, \texttt{PoseStamped}, \texttt{Clock}, \texttt{Imu}) to \texttt{meta/interfaces.txt}:
\begin{cmdbox}
printf '%s\n' \
  geometry_msgs/msg/PoseStamped \
  rosgraph_msgs/msg/Clock \
  sensor_msgs/msg/Imu \
  geometry_msgs/msg/Pose2D \
  >> ~/ros2_ws/src/sim_ros2_interface/meta/interfaces.txt
\end{cmdbox}

  \item Build cleanly inside the Pixi environment (this automatically links the plugin and Lua scripts into CoppeliaSim):
\begin{cmdbox}
cd ~/ros2_ws
export COPPELIASIM_ROOT_DIR="/Applications/coppeliaSim.app/Contents/Resources"
pixi run --manifest-path ~/ros_jazzy_pixi/pixi.toml \
  colcon build --symlink-install --cmake-clean-cache
\end{cmdbox}
\end{enumerate}

% =============================================================================
\section{Wrap CoppeliaSim so Applications/Finder Launches Inherit ROS 2}
% =============================================================================

macOS GUI apps launched from Applications, Finder, or Spotlight do not inherit your terminal's \texttt{AMENT\_PREFIX\_PATH}. Wrapping the launcher (and removing the old bundle signature) lets you open \textbf{CoppeliaSim} directly from Applications:

\begin{enumerate}[leftmargin=1.5em]
  \item Rename the original binary to \texttt{coppeliaSim\_bin}:
\begin{cmdbox}
mv /Applications/coppeliaSim.app/Contents/MacOS/coppeliaSim \
  /Applications/coppeliaSim.app/Contents/MacOS/coppeliaSim_bin
\end{cmdbox}

  \item Create the ROS 2 wrapper script at \texttt{/Applications/coppeliaSim.app/Contents/MacOS/coppeliaSim}:
\begin{cmdbox}
printf '%s\n' \
  '#!/usr/bin/env bash' \
  'if [ -z "$AMENT_PREFIX_PATH" ]; then' \
  '    MANIFEST="$HOME/ros_jazzy_pixi/pixi.toml"' \
  '    hook="$("$HOME/.pixi/bin/pixi" shell-hook -s bash --manifest-path "$MANIFEST" 2>/dev/null)"' \
  '    eval "$hook"' \
  'fi' \
  'export ROS_AUTOMATIC_DISCOVERY_RANGE=LOCALHOST' \
  'export ROS_DOMAIN_ID=42' \
  'export COPPELIASIM_ROOT_DIR="/Applications/coppeliaSim.app/Contents/Resources"' \
  'exec "$(dirname "$0")/coppeliaSim_bin" "$@"' \
  > /Applications/coppeliaSim.app/Contents/MacOS/coppeliaSim
chmod +x /Applications/coppeliaSim.app/Contents/MacOS/coppeliaSim
\end{cmdbox}

  \item Remove the stale code signature and quarantine flags so macOS does not report the modified app as damaged:
\begin{cmdbox}
rm -rf /Applications/coppeliaSim.app/Contents/_CodeSignature
xattr -cr /Applications/coppeliaSim.app
\end{cmdbox}
\end{enumerate}

% =============================================================================
\section{Verify Everything Together}
% =============================================================================

\begin{enumerate}[leftmargin=1.5em]
  \item In CoppeliaSim, open a ROS 2 lab scene and press \textbf{Play}.
  \item Confirm there are no \texttt{simROS2} or \texttt{plugin 'ROS2': load failed} errors in the CoppeliaSim console.
  \item In a second terminal, run \texttt{cd \textasciitilde/ros\_jazzy\_pixi \&\& pixi run ros2 topic list} and confirm the simulation topics appear.
\end{enumerate}

% =============================================================================
\section{Daily Workflow (Quick Reference)}
% =============================================================================

Once setup is complete, you never need to rebuild unless you edit \texttt{sim\_ros2\_interface}. Whenever you want to work on a lab, use two terminals:

\begin{itemize}[leftmargin=1.5em]
  \item \textbf{Applications / Terminal 1 (Start CoppeliaSim with ROS 2):}
  Open \textbf{CoppeliaSim} directly from Applications/Spotlight, or run in Terminal:
\begin{cmdbox}
/Applications/coppeliaSim.app/Contents/MacOS/coppeliaSim
\end{cmdbox}
  \item \textbf{Terminal 2 (Run ROS 2 nodes, topics, and tools):}
\begin{cmdbox}
cd ~/ros_jazzy_pixi
pixi shell
\end{cmdbox}
  Inside that shell, run any ROS 2 command (e.g., \texttt{ros2 topic list}, \texttt{ros2 topic echo /odom}, \texttt{ros2 run teleop\_twist\_keyboard teleop\_twist\_keyboard}, or \texttt{plotjuggler}).
\end{itemize}

\end{document}
